
% BEGIN bib
 
\par\medskip\noindent\textbf{References}\par
  
\par\noindent  [1]  B. Alarcón, R. Gutiérrez, and S. Lucas.  Context-sensitive dependency pairs.  IC, 208(8):922–968, 2010.  doi:10.1016/J.IC.2010.03.003.  
 
\par\noindent  [2]  T. Aoto and T. Yamada.  Argument filterings and usable rules for simply typed dependency pairs.  In S. Ghilardi and R. Sebastiani, editors, Proc. FroCoS, pages 117–132, 2009.  doi:10.1007/978-3-642-04222-5\_7.  
 
\par\noindent  [3]  T. Arts and J. Giesl.  Termination of term rewriting using dependency pairs.  TCS, 236(1–2):133–178, 2000.  doi:10.1016/S0304-3975(99)00207-8.  
 
\par\noindent  [4]  D. Beyer, A. Cimatti, A. Griggio, M. E. Keremoglu, and R. Sebastiani.  Software model checking via large-block encoding.  In A. Biere and C. Pixley, editors, Proc. FMCAD, pages 25–32, 2009.  doi:10.1109/FMCAD.2009.5351147.  
 
\par\noindent  [5]  M. Codish, C. Fuhs, J. Giesl, and P. Schneider-Kamp.  Lazy abstraction for size-change termination.  In C. G. Fermüller and A. Voronkov, editors, Proc. LPAR (Yogyakarta), pages 217–232, 2010.  doi:10.1007/978-3-642-16242-8\_16.  
 
\par\noindent  [6]  Community.  Termination competition (TermCOMP).  URL: https://termination-portal.org/wiki/Termination\_Competition.  
 
\par\noindent  [7]  Community.  The termination problem database (TPDB).  URL: https://github.com/TermCOMP/TPDB.  
 
\par\noindent  [8]  L. de Moura and N. Bjørner.  Z3: An efficient SMT solver.  In C.R. Ramakrishnan and J. Rehof, editors, Proc. TACAS, pages 337–340, 2008.  doi:10.1007/978-3-540-78800-3\_24.  
 
\par\noindent  [9]  N. Dershowitz.  Hierarchical termination.  In N. Dershowitz and N. Lindenstrauss, editors, Proc. CTRS, pages 89–105, 1995.  doi:10.1007/3-540-60381-6\_6.  
 
\par\noindent  [10]  S. Falke and D. Kapur.  A term rewriting approach to the automated termination analysis of imperative programs.  In R. A. Schmidt, editor, Proc. CADE, pages 277–293, 2009.  doi:10.1007/978-3-642-02959-2\_22.  
 
\par\noindent  [11]  S. Falke, D. Kapur, and C. Sinz.  Termination analysis of C programs using compiler intermediate languages.  In M. Schmidt-Schauß, editor, Proc. RTA, pages 41–50, 2011.  doi:10.4230/LIPIcs.RTA.2011.41.  
 
\par\noindent  [12]  C. Fuhs, J. Giesl, M. Plücker, P. Schneider-Kamp, and S. Falke.  Proving termination of integer term rewriting.  In R. Treinen, editor, Proc. RTA, pages 32–47, 2009.  doi:10.1007/978-3-642-02348-4\_3.  
 
\par\noindent  [13]  C. Fuhs and C. Kop.  A static higher-order dependency pair framework.  In L. Caires, editor, Proc. ESOP, pages 752–782, 2019.  doi:10.1007/978-3-030-17184-1\_27.  
 
\par\noindent  [14]  C. Fuhs, C. Kop, and N. Nishida.  Verifying procedural programs via constrained rewriting induction.  ACM TOCL, 18(2):14:1–14:50, 2017.  doi:10.1145/3060143.  
 
\par\noindent  [15]  J. Giesl, C. Aschermann, M. Brockschmidt, F. Emmes, F. Frohn, C. Fuhs, J. Hensel, C. Otto, M. Plücker, P. Schneider-Kamp, T. Ströder, S. Swiderski, and R. Thiemann.  Analyzing program termination and complexity automatically with AProVE.  JAR, 58(1):3–31, 2017.  doi:10.1007/s10817-016-9388-y.  
 
\par\noindent  [16]  J. Giesl, M. Raffelsieper, P. Schneider-Kamp, S. Swiderski, and R. Thiemann.  Automated termination proofs for Haskell by term rewriting.  ACM TOPLAS, 33(2):7:1–7:39, 2011.  doi:10.1145/1890028.1890030.  
 
\par\noindent  [17]  J. Giesl, T. Ströder, P. Schneider-Kamp, F. Emmes, and C. Fuhs.  Symbolic evaluation graphs and term rewriting: a general methodology for analyzing logic programs.  In A. King, editor, Proc. PPDP, pages 1–12, 2012.  doi:10.1145/2370776.2370778.  
 
\par\noindent  [18]  J. Giesl, R. Thiemann, and P. Schneider-Kamp.  The dependency pair framework: combining techniques for automated termination proofs.  In F. Baader and A. Voronkov, editors, Proc. LPAR, pages 301–331, 2005.  doi:10.1007/978-3-540-32275-7\_21.  
 
\par\noindent  [19]  J. Giesl, R. Thiemann, P. Schneider-Kamp, and S. Falke.  Mechanizing and improving dependency pairs.  JAR, 37(3):155–203, 2006.  doi:10.1007/s10817-006-9057-7.  
 
\par\noindent  [20]  L. Guo, K. Hagens, C. Kop, and D. Vale.  Higher-order constrained dependency pairs for (universal) computability.  pre-publication Arxiv copy of [21] with additional appendix.  doi:10.48550/arXiv.2406.19379.  
 
\par\noindent  [21]  L. Guo, K. Hagens, C. Kop, and D. Vale.  Higher-order constrained dependency pairs for (universal) computability.  In R. Královič and A. Kučera, editors, Proc. MFCS, pages 57:1–57:15, 2024.  doi:10.4230/LIPIcs.MFCS.2024.57.  
 
\par\noindent  [22]  L. Guo and C. Kop.  Higher-order LCTRSs and their termination.  In S. Weirich, editor, Proc. ESOP, pages 331–357, 2024.  doi:10.1007/978-3-031-57267-8\_13.  
 
\par\noindent  [23]  N. Hirokawa and A. Middeldorp.  Automating the dependency pair method.  IC, 199(1-2):172–199, 2005.  doi:10.1016/J.IC.2004.10.004.  
 
\par\noindent  [24]  N. Hirokawa and A. Middeldorp.  Tyrolean Termination Tool: techniques and features.  IC, 205(4):474–511, 2007.  doi:10.1016/j.ic.2006.08.010.  
 
\par\noindent  [25]  C. Kop.  WANDA – a higher order termination tool (system description).  In Z. M. Ariola, editor, Proc. FSCD, pages 36:1–36:19, 2020.  doi:10.4230/LIPICS.FSCD.2020.36.  
 
\par\noindent  [26]  C. Kop.  Cutting a proof into bite-sized chunks: incrementally proving termination in higher-order term rewriting.  In A. P. Felty, editor, Proc. FSCD, pages 1:1–1:17, 2022.  doi:10.4230/LIPIcs.FSCD.2022.1.  
 
\par\noindent  [27]  C. Kop et al.  The Cora analyzer.  URL: https://github.com/hezzel/cora.  
 
\par\noindent  [28]  C. Kop et al.  hezzel/cora: FSCD 2025.  doi:10.5281/zenodo.15318964.  
 
\par\noindent  [29]  C. Kop and N. Nishida.  Term rewriting with logical constraints.  In P. Fontaine, C. Ringeissen, and R. A. Schmidt, editors, Proc. FroCoS, pages 343–358, 2013.  doi:10.1007/978-3-642-40885-4\_24.  
 
\par\noindent  [30]  C. Kop and F. van Raamsdonk.  Dynamic dependency pairs for algebraic functional systems.  LMCS, 8(2), 2012.  doi:10.2168/LMCS-8(2:10)2012.  
 
\par\noindent  [31]  M. R. K. Krishna Rao.  Completeness of hierarchical combinations of term rewriting systems.  In R. K. Shyamasundar, editor, Proc. FSTTCS, pages 125–138, 1993.  doi:10.1007/3-540-57529-4\_48.  
 
\par\noindent  [32]  M. R. K. Krishna Rao.  Simple termination of hierarchical combinations of term rewriting systems.  In M. Hagiya and J. C. Mitchell, editors, Proc. TACS, pages 203–223, 1994.  doi:10.1007/3-540-57887-0\_97.  
 
\par\noindent  [33]  M. R. K. Krishna Rao.  Semi-completeness of hierarchical and super-hierarchical combinations of term rewriting systems.  In P. D. Mosses, M. Nielsen, and M. I. Schwartzbach, editors, Proc. CAAP, pages 379–393, 1995.  doi:10.1007/3-540-59293-8\_208.  
 
\par\noindent  [34]  K. Kusakari, Y. Isogai, M. Sakai, and F. Blanqui.  Static dependency pair method based on strong computability for higher-order rewrite systems.  IEICE Trans. Inf. Syst., E92.D(10):2007–2015, 2009.  doi:10.1587/transinf.E92.D.2007.  
 
\par\noindent  [35]  K. Kusakari and M. Sakai.  Enhancing dependency pair method using strong computability in simply-typed term rewriting.  AAECC, 18(5):407–431, 2007.  doi:10.1007/s00200-007-0046-9.  
 
\par\noindent  [36]  K. Kusakari and M. Sakai.  Static dependency pair method for simply-typed term rewriting and related techniques.  IEICE Trans. Inf. Syst., E92.D(2):235–247, 2009.  doi:10.1587/transinf.E92.D.235.  
 
\par\noindent  [37]  C. S. Lee, N. D. Jones, and A. M. Ben-Amram.  The size-change principle for program termination.  In C. Hankin and D. Schmidt, editors, Proc. POPL, pages 81–92, 2001.  doi:10.1145/360204.360210.  
 
\par\noindent  [38]  S. Lucas.  Context-sensitive computations in functional and functional logic programs.  JFLP, 1998(1), 1998.  URL: http://danae.uni-muenster.de/lehre/kuchen/JFLP/articles/1998/A98-01/A98-01.html.  
 
\par\noindent  [39]  C. Otto, M. Brockschmidt, C. von Essen, and J. Giesl.  Automated termination analysis of Java bytecode by term rewriting.  In C. Lynch, editor, Proc. RTA, pages 259–276, 2010.  doi:10.4230/LIPIcs.RTA.2010.259.  
 
\par\noindent  [40]  S. Suzuki, K. Kusakari, and F. Blanqui.  Argument filterings and usable rules in higher-order rewrite systems.  IPSJ Online Trans., 4:114–125, 2011.  doi:10.2197/ipsjtrans.4.114.  
 
\par\noindent  [41]  R. Thiemann.  The DP framework for proving termination of term rewriting.  PhD thesis, RWTH Aachen University, Germany, 2007.  URL: http://darwin.bth.rwth-aachen.de/opus3/volltexte/2007/2066/.  
 
\par\noindent  [42]  R. Thiemann and J. Giesl.  The size-change principle and dependency pairs for termination of term rewriting.  AAECC, 16(4):229–270, 2005.  doi:10.1007/S00200-005-0179-7.  
 
\par\noindent  [43]  J. van de Pol.  Termination proofs for higher-order rewrite systems.  In J. Heering, K. Meinke, B. Möller, and T. Nipkow, editors, Proc. HOA, pages 305–325, 1993.  doi:10.1007/3-540-58233-9\_14.  
 
\par\noindent  [44]  F. van Raamsdonk.  Translating logic programs into conditional rewriting systems.  In L. Naish, editor, Proc. ICLP, pages 168–182, 1997.  doi:10.7551/mitpress/4299.003.0018.  
 
\par\noindent  [45]  H. Zantema.  Termination of term rewriting by semantic labelling.  FI, 24(1/2):89–105, 1995.  doi:10.3233/FI-1995-24124.  
  
% END bib
